\subsection{Littlewood--Paley 理论}
\label{sec:ch8-littlewood-paley-theory}

Plancherel 定理表明, 把 $L^2$ 函数的 Fourier 变换乘以模为 $1$ 的函数, 所得仍是 $L^2$ 函数. 类似地, 把 $L^2$ 函数 Fourier 级数的各项乘以 $\pm1$, 所得仍是 $L^2$ 函数. 但对 $L^p$, $p\ne2$, 这两项性质都不成立. 因而一个函数是否属于 $L^p$, 不只取决于其 Fourier 变换或 Fourier 系数的大小.

Littlewood--Paley 理论在 $L^p$ 中为 Plancherel 定理的结论提供了部分替代. 它表明, 若在二进块内用 $\pm1$ 修改 Fourier 变换或 Fourier 系数, 则 $L^p$ 归属关系不变. 例如, 对 $L^p$ 函数的 Fourier 级数, 若对所有指标位于 $2^k$ 与 $2^{k+1}$ 之间的系数赋予同一因子 $1$ 或 $-1$, 所得仍是另一个 $L^p$ 函数的 Fourier 级数.

\hypertarget{chapter-08-page-172}{}
先对定义在 $\mathbb R$ 上的函数证明 Fourier 变换的相应结果. 令
\[
 \Delta_j=(-2^{j+1},-2^j]\cup[2^j,2^{j+1}),
\]
并定义
\[
 \widehat{S_jf}(\xi)=\chi_{\Delta_j}(\xi)\widehat f(\xi),
 \qquad j\in\mathbb Z.
\]
区间 $[2^j,2^{j+1})$ 称为二进区间. 若 $f\in L^2$, 则 Plancherel 定理给出
\begin{equation}
\label{eq:ch8-lp-l2-identity}
 \left\|\left(\sum_j|S_jf|^2\right)^{1/2}\right\|_2=\|f\|_2.
\end{equation}
Littlewood--Paley 定理说, 这两个量在 $L^p$ 中也可比较.

\begin{Theorem}[Littlewood--Paley]
\label{thm:ch8-littlewood-paley}
设 $f\in L^p(\mathbb R)$, $1<p<\infty$. 则存在正常数 $c_p,C_p$, 使
\begin{equation}
\label{eq:ch8-littlewood-paley}
 c_p\|f\|_p
 \le\left\|\left(\sum_j|S_jf|^2\right)^{1/2}\right\|_p
 \le C_p\|f\|_p.
\end{equation}
\end{Theorem}

将把定理~\ref{thm:ch8-littlewood-paley} 作为一个类似结果的推论加以证明, 但后者用光滑函数而不是区间示性函数定义算子. 取非负 $\psi\in\mathcal S(\mathbb R)$, 使其支集包含于 $1/2\le|\xi|\le4$, 并在 $1\le|\xi|\le2$ 上等于 $1$. 定义
\[
 \psi_j(\xi)=\psi(2^{-j}\xi),
 \qquad
 \widehat{\widetilde S_jf}(\xi)=\psi_j(\xi)\widehat f(\xi).
\]
于是立即有 $\widetilde S_jS_j=S_j$.

\begin{Theorem}
\label{thm:ch8-smooth-square-function}
设 $f\in L^p(\mathbb R)$, $1<p<\infty$. 则存在常数 $C_p$, 使
\begin{equation}
\label{eq:ch8-smooth-square-function}
 \left\|\left(\sum_j|\widetilde S_jf|^2\right)^{1/2}\right\|_p
 \le C_p\|f\|_p.
\end{equation}
\end{Theorem}

\begin{Proof}
记 $\check\psi$ 为 $\psi$ 的 Fourier 逆变换, $\check\psi_j(x)=2^j\check\psi(2^jx)$. 则 $\widehat{\check\psi_j}=\psi_j$ 且 $\widetilde S_jf=\check\psi_j*f$. 因而只需证明把 $f$ 映为序列 $\{\widetilde S_jf\}$ 的向量值算子从 $L^p$ 到 $L^p(\ell^2)$ 有界. 当 $p=2$ 时, Plancherel 定理给出
\[
 \left\|\left(\sum_j|\widetilde S_jf|^2\right)^{1/2}\right\|_2^2
 =\int_{\mathbb R}\sum_j|\psi_j(\xi)|^2|\widehat f(\xi)|^2\,d\xi
 \le3\|f\|_2^2,
\]
因为对每个 $\xi$, 至多三个 $\psi_j$ 非零.

\hypertarget{chapter-08-page-173}{}
为证明 \eqref{eq:ch8-smooth-square-function}, 需要验证向量值算子的核 $\{\check\psi_j\}$ 满足 Hörmander 条件. 只需证明
\[
 \|\{(\check\psi_j)'(x)\}\|_{\ell^2}\le C|x|^{-2}.
\]
这也可由命题~\ref{prop:ch5-dini-implies-hormander} 的向量值版本推出. 事实上,
\[
 \left(\sum_j|(\check\psi_j)'(x)|^2\right)^{1/2}
 \le\sum_j2^{2j}|(\check\psi)'(2^jx)|.
\]
又因 $\check\psi\in\mathcal S$, 有 $|(\check\psi)'(x)|\le C\min(1,|x|^{-3})$. 固定 $i$ 使 $2^{-i}\le|x|<2^{-i+1}$. 则上式不超过
\[
 C\sum_{j\le i}2^{2j}+C|x|^{-3}\sum_{j>i}2^{-j}
 \le C|x|^{-2}.
\]
\end{Proof}

\begin{Proof}[定理~\ref{thm:ch8-littlewood-paley} 的证明]
由推论~\ref{cor:ch8-interval-projections-vector-valued} 和恒等式 $\widetilde S_jS_j=S_j$,
\[
 \left\|\left(\sum_j|S_jf|^2\right)^{1/2}\right\|_p
 =\left\|\left(\sum_j|S_j\widetilde S_jf|^2\right)^{1/2}\right\|_p
 \le C\left\|\left(\sum_j|\widetilde S_jf|^2\right)^{1/2}\right\|_p.
\]
所以定理~\ref{thm:ch8-smooth-square-function} 立即给出 \eqref{eq:ch8-littlewood-paley} 的右端不等式.

由极化恒等式和 \eqref{eq:ch8-lp-l2-identity},
\[
 \int_{\mathbb R}\sum_jS_jf\,\overline{S_jg}
 =\int_{\mathbb R}f\overline g.
\]
因此
\[
\begin{aligned}
 \|f\|_p
 &=\sup_{\|g\|_{p'}\le1}\left|\int_{\mathbb R}f\overline g\right|\\
 &=\sup_{\|g\|_{p'}\le1}\left|\int_{\mathbb R}\sum_jS_jf\,\overline{S_jg}\right|\\
 &\le\sup_{\|g\|_{p'}\le1}
 \left\|\left(\sum_j|S_jf|^2\right)^{1/2}\right\|_p
 \left\|\left(\sum_j|S_jg|^2\right)^{1/2}\right\|_{p'}\\
 &\le C_p\left\|\left(\sum_j|S_jf|^2\right)^{1/2}\right\|_p.
\end{aligned}
\]
\end{Proof}

这些结果有两种不同的 $\mathbb R^n$ 推广. 定理~\ref{thm:ch8-smooth-square-function} 的一个推广用支于环带 $2^{j-1}<|\xi|<2^{j+1}$ 的光滑函数分割 Fourier 变换; 定理~\ref{thm:ch8-littlewood-paley} 的另一推广则用坐标轴上二进区间乘积的示性函数作分割.

\hypertarget{chapter-08-page-174}{}
示性函数环带的推广不成立, 这一点由下文定理~\ref{thm:ch8-ball-multiplier-negative} 可知.

\begin{Theorem}
\label{thm:ch8-smooth-square-rn}
给定 $\psi\in\mathcal S(\mathbb R^n)$ 且 $\psi(0)=0$, 对 $j\in\mathbb Z$ 定义
\[
 \widehat{S_jf}(\xi)=\psi(2^{-j}\xi)\widehat f(\xi).
\]
则当 $1<p<\infty$ 时,
\begin{equation}
\label{eq:ch8-smooth-square-rn}
 \left\|\left(\sum_j|S_jf|^2\right)^{1/2}\right\|_p
 \le C_p\|f\|_p.
\end{equation}
此外, 若对所有 $\xi\ne0$,
\begin{equation}
\label{eq:ch8-smooth-partition}
 \sum_j|\psi(2^{-j}\xi)|^2=C,
\end{equation}
则还有
\begin{equation}
\label{eq:ch8-smooth-square-rn-lower}
 \|f\|_p\le C_p
 \left\|\left(\sum_j|S_jf|^2\right)^{1/2}\right\|_p.
\end{equation}
\end{Theorem}

由于 $\psi\in\mathcal S$ 且 $\psi(0)=0$, 有 $\sum_j|\psi(2^{-j}\xi)|^2<C$. 因此第一部分的证明与定理~\ref{thm:ch8-smooth-square-function} 完全相同; 第二部分也与定理~\ref{thm:ch8-littlewood-paley} 相应部分相同, 因为 \eqref{eq:ch8-smooth-partition} 给出 $p=2$ 时的等式, 至多相差常数.

可以简单地构造满足 \eqref{eq:ch8-smooth-partition} 的函数. 固定非负、径向递减的 $\phi\in\mathcal S(\mathbb R^n)$, 使 $|\xi|\le1/2$ 时 $\phi(\xi)=1$, $|\xi|\ge1$ 时 $\phi(\xi)=0$. 定义
\[
 \psi(\xi)^2=\phi(\xi/2)-\phi(\xi).
\]
则立即有
\[
 \sum_j|\psi(2^{-j}\xi)|^2=1,
 \qquad\xi\ne0.
\]

下一结果只在 $\mathbb R^2$ 上陈述和证明, 高维推广可立即由归纳得到. 平面中的二进矩形是实直线上二进区间的 Cartesian 乘积. 相应算子是 $S_j^1S_k^2$, 其中
\[
 \widehat{S_j^1f}(\xi_1,\xi_2)=\chi_{\Delta_j}(\xi_1)\widehat f(\xi_1,\xi_2),
 \qquad
 \widehat{S_k^2f}(\xi_1,\xi_2)=\chi_{\Delta_k}(\xi_2)\widehat f(\xi_1,\xi_2).
\]

\hypertarget{chapter-08-page-175}{}
\begin{Theorem}
\label{thm:ch8-product-littlewood-paley}
设 $f\in L^p(\mathbb R^2)$, $1<p<\infty$. 则存在正常数 $c_p,C_p$, 使
\[
 c_p\|f\|_p
 \le\left\|\left(\sum_{j,k}|S_j^1S_k^2f|^2\right)^{1/2}\right\|_p
 \le C_p\|f\|_p.
\]
\end{Theorem}

\begin{Proof}
用定理~\ref{thm:ch8-smooth-square-function} 证明中的同一记号和论证,
\begin{equation}
\label{eq:ch8-smooth-sequence-vector}
 \left\|\left(\sum_j|\widetilde S_jf_j|^2\right)^{1/2}\right\|_p
 \le C_p\left\|\left(\sum_j|f_j|^2\right)^{1/2}\right\|_p.
\end{equation}
由 Fubini 定理, 即使 $f_j$ 是二变量函数而 $\widetilde S_j$ 只作用于其中一个变量, \eqref{eq:ch8-smooth-sequence-vector} 仍成立. 再按定理~\ref{thm:ch8-littlewood-paley} 的证明, 推论~\ref{cor:ch8-interval-projections-vector-valued} 可直接用于此情形, 因而还可用 $S_j$ 替换 $\widetilde S_j$, 因为 $\widetilde S_jS_j=S_j$. 先应用 \eqref{eq:ch8-smooth-sequence-vector}, 再应用定理~\ref{thm:ch8-littlewood-paley}, 并使用 Fubini 定理, 得
\[
 \left\|\left(\sum_{j,k}|S_j^1S_k^2f|^2\right)^{1/2}\right\|_p
 \le C_p\left\|\left(\sum_k|S_k^2f|^2\right)^{1/2}\right\|_p
 \le C_p^2\|f\|_p.
\]
反向不等式由同一对偶论证得到.
\end{Proof}
